%
% 第6章：跳跃连接与注意力机制

\clearpage
\cleardoublepage

\chapter{跳跃连接与注意力机制}

\section{跳跃连接}

跳跃连接（捷径）绕过层以实现深度网络中的梯度流动。

\subsection{残差连接（ResNet）}

残差学习框架：
\begin{equation}
\mathbf{y} = \mathcal{F}(\mathbf{x}) + \mathbf{x}
\end{equation}
其中 $\mathcal{F}$ 是学习的残差映射。

属性：
\begin{itemize}
    \item \textbf{实现极深网络（100+层）的训练}
    \item \textbf{缓解梯度消失/爆炸问题}
    \item \textbf{减少退化问题}
    \item \textbf{易于实现，计算开销小}
\end{itemize}

如果维度匹配：$\mathbf{y} = \mathcal{F}(\mathbf{x}, \{W_i\}) + \mathbf{x}$

如果维度不同：$\mathbf{y} = \mathcal{F}(\mathbf{x}, \{W_i\}) + W_s\mathbf{x}$

\section{注意力机制}

注意力允许模型关注输入的相关部分。

\subsection{自注意力}

注意力函数：
\begin{equation}
\text{Attention}(\mathbf{Q}, \mathbf{K}, \mathbf{V}) = \text{softmax}\left(\frac{\mathbf{Q}\mathbf{K}^\top}{\sqrt{d_k}}\right)\mathbf{V}
\end{equation}

其中：
\begin{itemize}
    \item $\mathbf{Q} = \mathbf{X}\mathbf{W}_Q$（查询）
    \item $\mathbf{K} = \mathbf{X}\mathbf{W}_K$（键）
    \item $\mathbf{V} = \mathbf{X}\mathbf{W}_V$（值）
\end{itemize}

缩放因子 $\sqrt{d_k}$ 防止大 $d_k$ 时的梯度消失。

自注意力属性：
\begin{itemize}
    \item \textbf{允许建模长期依赖}
    \item \textbf{计算 $O(n^2)$ 但可并行化}
    \item \textbf{可解释：} 注意力权重显示相关性
\end{itemize}

\section{本章小结}

\begin{enumerate}
    \item 跳跃连接通过改善梯度流动实现深度网络
    \item 残差连接是使用最广泛的跳跃连接
    \item 注意力机制允许动态关注相关特征
    \item 自注意力是Transformer的基础
    \item 通道注意力（SE）提供高效的特征重校准
\end{enumerate}
